\section{Special cases, comparison with the \texorpdfstring{\(\pi\)}{pi}
proof, and limitations}\label{sec:consequences}

\subsection{Rational powers and rational affine changes}

\begin{lemma}\label{lem:affine}
If \(b\ge2\), \(\ELB(x,b)\), and \(a,c\in\Q\) with \(a\ne0\),
then \(\ELB(ax+c,b)\).
\end{lemma}
\begin{proof}
For a fixed \(\nu>b\), choose \(b<\nu'<\nu\).
Write \(a=A/B\), \(c=C/D\), with \(B,D>0\).
For any \(p/q\), the rational number \((p/q-c)/a\) has an integer
representation with denominator \(|A|Dq\).
Consequently, for sufficiently large \(q\),
\[
 |ax+c-p/q|\ge |a|(|A|Dq)^{-\nu'}.
\]
The fixed positive factor
\(|a|(|A|D)^{-\nu'}\) is eventually at least
\(q^{-(\nu-\nu')}\), proving the assertion.
\end{proof}

\begin{corollary}\label{cor:rational-powers}
Suppose \(\alpha>0\), \(n\in\N_{>0}\), and
\(\alpha^n=r\in\Q_{>0}\setminus\{1\}\).
Then
\[
 \mu(\log\alpha)=2.
\]
In particular,
\(\mu(\log\sqrt2)=2\).
\end{corollary}
\begin{proof}
The identity \(\log\alpha=(1/n)\log r\), together with
Theorem~\ref{thm:rational}, Lemma~\ref{lem:affine}, and
Lemma~\ref{lem:endpoint}, proves the result.
\end{proof}

This sharp conclusion is stronger than the degree bound in this
particular family. The proof uses rational scaling, not invariance
under arbitrary algebraic scaling.

\subsection{Arctangents of rational numbers}

\begin{corollary}\label{cor:arctangent}
For every \(r\in\Q\setminus\{0\}\),
\[
 \mu(\arctan r)=2.
\]
\end{corollary}
\begin{proof}
Use \(F=\Q(i)\), \(\gamma=2i\),
\(\xi=(1+ir)/(1-ir)\), and \(x=\arctan r\ne0\).
The identity
\[
 e^{2i\arctan r}=\frac{1+ir}{1-ir}
\]
establishes compatibility for the given embedding. Its complex
conjugate establishes compatibility for the other embedding.
Thus \(s=d=2\).
For \(r=1\), \(\xi=i\) is a root of unity; the root-of-unity
branch of the theorem is genuinely needed here.
\end{proof}

\subsection{Quadratic-normalized logarithms}

\begin{corollary}\label{cor:quadratic}
Let \(a,b,D\in\Q\), \(D>0\), and suppose
\[
 a^2-b^2D=1,\qquad \alpha=a+b\sqrt D>1.
\]
Then
\[
 \mu\left(\frac{\log(a+b\sqrt D)}{\sqrt D}\right)=2.
\]
In particular,
\begin{equation}\label{eq:quadratic-example}
 \mu\left(\frac{\log(3+2\sqrt2)}{\sqrt2}\right)=2.
\end{equation}
\end{corollary}
\begin{proof}
Take \(F=\Q(\sqrt D)\), \(\gamma=\sqrt D\),
\(\xi=a+b\sqrt D\), and \(x=\log\alpha/\sqrt D\ne0\).
Every embedding sends \(\sqrt D\) to either \(\sqrt D\) or
\(-\sqrt D\). The first choice gives
\(e^{\sqrt D\,x}=\alpha\).
The norm identity gives \(a-b\sqrt D=\alpha^{-1}\), so the
second choice gives
\[
 e^{-\sqrt D\,x}=\alpha^{-1}=a-b\sqrt D.
\]
Every embedding is compatible, whether the field has degree one or two.
\end{proof}

The normalization in this corollary is essential to the asserted
specialization. It does not imply
\(\mu(\log(3+2\sqrt2))=2\) by removing the factor \(\sqrt2\).
The general algebraic-base result gives the separate bound
\(2\le\mu(\log(3+2\sqrt2))\le4\).

\subsection{How the argument differs from the original
\texorpdfstring{\(\pi\)}{pi} proof}

The relation to \cite{OAIpi} is at the level of proved machinery,
not a substitution into its final theorem.
The following comparison fixes the main changes.

{\small
\begin{longtable}{@{}>{\raggedright\arraybackslash}p{.22\textwidth}
>{\raggedright\arraybackslash}p{.34\textwidth}
>{\raggedright\arraybackslash}p{.36\textwidth}@{}}
\toprule
Feature & Original \(\pi\) construction & Logarithmic extension\\
\midrule
\endfirsthead
\toprule
Feature & Original \(\pi\) construction & Logarithmic extension\\
\midrule
\endhead
\bottomrule
\endfoot
Exponential identity
& \(e^{2i\pi}=1\)
& \(e^{\log r}=r\), or \(e^{\sigma(\gamma)x}=\sigma(\xi)\)\\[5pt]
Multiplicative centers
& Common fiber \(Y=1\)
& Distinct powers \(r^j\) or \(\xi^j\); common fiber after
root-of-unity reduction\\[5pt]
Local coordinates
& \(t=Y-1\)
& \(t=Y/y_j-1\); the same logarithmic frame\\[5pt]
Arithmetic
& Gaussian-integer norm after denominator clearing
& Rational integers for rational logarithms; an algebraic integer
and its full field norm in general\\[5pt]
Analytic saving
& Repeated \((a,d)\) coefficient rows
& The same repetition argument at every compatible embedding;
direct growth bounds at the others\\[5pt]
Final coefficient
& \(1\) in the normalized arithmetic comparison
& \(d/s\), giving threshold \(2d/s\)\\
\end{longtable}
}

For \(\pi\), the compatible data may be viewed as
\(F=\Q(i)\), \(\gamma=2i\), \(\xi=1\), \(x=\pi\):
both embeddings are compatible. For \(\log r\), the field is
\(\Q\), the slope is \(1\), and the real centers require the
distinct-fiber argument. The analytic translation is preserved
exactly, while the radius constants change.
The proof of the new interpolation theorem, the identification with
\eqref{eq:matrix}, the field-denominator estimates, and the estimates
at all embeddings are indispensable additions.

\subsection{Why no theorem for all positive real bases is possible}

\begin{proposition}\label{prop:counterexamples}
The map \(\alpha\mapsto\log\alpha\) on positive real numbers is
surjective onto \(\R\). There is a positive real \(\alpha\ne1\)
with rational logarithm, and there is such an \(\alpha\) whose
irrational logarithm has no finite irrationality exponent.
\end{proposition}
\begin{proof}
For every \(y\in\R\), \(\alpha=e^y>0\) satisfies
\(\log\alpha=y\). Taking \(y=1\) gives the rational example.
For the second, let
\[
 L=\sum_{n=0}^\infty10^{-n!},\qquad \alpha=e^L.
\]
The standard Liouville argument applies: truncation through \(N\)
has denominator \(q_N=10^{N!}\), while its positive tail is at most
\(2q_N^{-(N+1)}\). These approximations prove that \(L\) is an
irrational Liouville number and admits no eventual lower bound of
any finite exponent. Thus \(\log\alpha=L\) has the claimed property.
\end{proof}

These examples concern arbitrary real inputs. They are not
counterexamples to the rational- or algebraic-base theorems.
The formal statements express the second conclusion as failure of
\(\ELB(\log\alpha,b)\) for every finite real \(b\), avoiding any
ambiguous convention for a real-valued supremum of an unbounded set.
